\section{Introduction}

AI coding agents increasingly operate as long-running processes with access to
shells, source trees, package managers, and external APIs. Project maintainers
steer these agents with instruction files such as \texttt{CLAUDE.md} and
\texttt{AGENTS.md}: do not push directly to \texttt{main}; run tests before
committing; never expose secrets; use only a sanctioned review tool before
destructive operations. These are not style preferences. They are behavioral
constraints that must hold across the agent's entire execution.

A cooperative agent will generally follow these constraints---but it forgets.
In long contexts, attention is diluted, earlier instructions drop out of the
effective window, and the agent takes a shortcut that violates a rule it
previously acknowledged. The agent is not adversarial; it is
\emph{cooperative but forgetful}. It would benefit from a mechanism that
makes its own declared constraints durable---surviving context loss, tool
switching, and subprocess indirection---and that reminds it when a constraint
is about to be violated.

\subsection{The Intent--Behavior Gap}

We distinguish three levels at which an agent's execution can be described,
following AgentSight's observation of a ``semantic gap'' between agent intent
and low-level actions~\cite{agentsight}.

\textbf{Intent.} What the agent declares it will and will not do: goals
(``fix this bug'') and self-constraints (``I must not leak secrets'').
Intent is expressed in natural language and, in ActPlane's design, formalized
in a compact DSL that the agent itself can write and maintain.

\textbf{Action.} The tool calls the agent issues through its runtime:
\texttt{read\_file}, \texttt{run\_command}, \texttt{edit\_file}. This is the
interface between the agent and the outside world, mediated by the agent
framework.

\textbf{Behavior.} The actual OS-level operations that result: \texttt{open},
\texttt{execve}, \texttt{connect}, \texttt{fork}. Behavior is where effects
really happen---files are read, processes are spawned, data reaches the
network.

The core problem is that the mapping from action to behavior is
\emph{many-to-many and opaque}. A single tool call
\texttt{run\_command("make")} can produce hundreds of system calls. The same
OS operation---\texttt{execve("git", ["push"])}---can be reached through a
direct tool call, a \texttt{bash~-c} string, a Python subprocess, or a
compiled helper binary. An action-level guard sees only the tool call; it
cannot predict or intercept the behavior it produces.

\subsection{Why Single-Level Enforcement Is Insufficient}

Enforcement can be placed at any of the three levels. Each alone is
insufficient.

\textbf{Intent-level} enforcement (prompt instructions, \texttt{CLAUDE.md})
is probabilistic: the agent may forget the constraint or override it in
reasoning.

\textbf{Action-level} enforcement (tool-call guards such as AgentSpec and
Progent~\cite{wang2025agentspec,shi2025progent}) is deterministic but
bypassable: it checks tool calls, not the OS behavior they produce. When the
agent routes an effect through a subprocess, an SDK, or a direct system
interface, the guard sees nothing.

\textbf{Behavior-level} enforcement (syscall filters, sandboxes, eBPF
enforcers) is unbypassable---every effect crosses the kernel boundary. But
isolated behavior-level enforcement returns only \texttt{-EPERM} or
\texttt{SIGKILL}; the agent receives no explanation and cannot recover.

The gap between intent and behavior is not bridged by any single level.
Intent-level constraints are not connected to behavior-level enforcement;
behavior-level enforcement is not connected back to intent-level feedback.

\subsection{ActPlane: Bridging Intent and Behavior}

ActPlane closes this gap in both directions.

\textbf{Downward: intent to behavior.} The agent declares behavioral
contracts in a DSL---voluntary confinement, analogous to \texttt{pledge()} in
OpenBSD, where a program restricts its own capabilities because it knows it
may be exploited. A coding agent can write and maintain this DSL just as it
edits any other project file. ActPlane compiles these declarations into
labeled information-flow rules and loads them into an eBPF/BPF-LSM backend.
Named labels propagate across process, file, and network objects at kernel
hooks, tracking which sources have influenced which subjects.

\textbf{Horizontal: across objects.} Labels encode execution history at the
behavior level. When a process reads a file labeled \texttt{SECRET}, the
process acquires that label. When it forks, the child inherits the label.
This cross-object propagation turns individual OS operations into checkable
behavioral state---enabling constraints like ``a process that has read
\texttt{.env} must not connect to external endpoints.''

\textbf{Upward: behavior to intent.} When a violation is detected, the
kernel reports the rule, the target, and the effect. ActPlane translates
this into corrective feedback delivered to the agent's context: which
constraint was violated, why the current operation conflicts with it, and
what alternative path (running a redactor, going through a review tool,
executing tests first) would satisfy the contract. The agent receives not
\texttt{Permission denied} but a reminder of its own declared intent.

The term \emph{agent harness} has come to mean everything around the
model---tool dispatch, memory, orchestration, guardrails, and
feedback~\cite{langchain-harness,fowler-harness}. Current harness
implementations place all of these components, including enforcement and
feedback, at the application layer. ActPlane extends the harness to the OS
level for the components that cannot work at the application layer alone:
behavioral enforcement and corrective feedback over the agent's actual
OS-level behavior. It is not a replacement for the application-layer harness
but a complement---the part that must reach below the tool API.

This design is not a new enforcement mechanism---CamQuery showed that
cross-channel label propagation and enforcement can be done in the
kernel~\cite{pasquier2018camquery}. Nor is corrective feedback for agents
new---AgentSpec demonstrated it at the tool-call
layer~\cite{wang2025agentspec}. And letting agents participate in their own
governance through information-flow control is already practiced at the
application layer by systems such as FIDES and
CaMeL~\cite{costa2025fides,debenedetti2025camel}. ActPlane's contribution is
bringing these capabilities together \emph{at the OS level}: cross-channel
labeled information-flow enforcement on the modern eBPF/BPF-LSM substrate,
driven by agent-declared behavioral contracts, with a feedback path that
connects kernel-detected violations back to the agent's intent---where
constraints survive context loss and cannot be bypassed below the tool API.

\subsection{Contributions}

This paper makes three contributions.

\begin{enumerate}
\item It defines the \emph{intent--behavior gap} for AI agents---the
disconnect between behavioral constraints declared at the intent level and
their enforcement at the OS level---and grounds it with a corpus study of
instruction files from 50 popular projects, showing that OS-enforceable
behavioral constraints are pervasive.

\item It presents ActPlane, a system that compiles agent-declared behavioral
contracts into eBPF/BPF-LSM labeled information-flow enforcement. The system
propagates labels across process, file, and network objects, checks
source-to-sink rules at kernel hooks, and returns corrective feedback that
reconnects behavior-level violations to intent-level remediation.

\item It evaluates ActPlane on bypass coverage across five tool paths, agent
recovery rate under four feedback conditions, false positives under benign
workloads, and per-event overhead, comparing against action-level and
behavior-level baselines.
\end{enumerate}
